% A linear iterative process for computing 6!, used in Section 1.2.1.
%
% The companion to factorial-recursive-process.tex, and the contrast between
% the two is the argument of the section. There the arrow bulges, because the
% process expands as deferred multiplications accumulate and contracts as they
% are performed. Here it is square: the process neither grows nor shrinks, and
% each line is the same size as the last -- three numbers.
%
% The arrow is drawn with right angles deliberately, for that reason.

\begin{tikzpicture}[
  x = 1mm, y = 4.6mm,
  line/.style = {anchor=west, font=\ttfamily\small, inner sep=0pt},
]

\node[line] (l1) at (0,  0) {factorial(6)};
\node[line] (l2) at (0, -1) {fact\_iter(1, 1, 6)};
\node[line]      at (0, -2) {fact\_iter(1, 2, 6)};
\node[line]      at (0, -3) {fact\_iter(2, 3, 6)};
\node[line]      at (0, -4) {fact\_iter(6, 4, 6)};
\node[line]      at (0, -5) {fact\_iter(24, 5, 6)};
\node[line]      at (0, -6) {fact\_iter(120, 6, 6)};
\node[line] (lw) at (0, -7) {fact\_iter(720, 7, 6)};
\node[line] (l9) at (0, -8) {720};

\draw[-{Triangle[open, length=1.8mm]}, thin]
  ([xshift=2mm] l1.east) -- ([xshift=8mm] l1.east);
\draw[thin]
  ([xshift=8mm] l1.east)
    -- ([xshift=10mm] lw.east |- l1.east)
    -- ([xshift=10mm] lw.east |- l9.east)
    -- ([xshift=8mm] l9.east);
\draw[-{Triangle[open, length=1.8mm]}, thin]
  ([xshift=8mm] l9.east) -- ([xshift=2mm] l9.east);

\end{tikzpicture}
